\chapter{Medium-Frequency Proprietary Shops}\label{ch:in:medium-frequency-proprietary-shops}

Most proprietary trading firms are small. At the end of 2024 the US broker-dealer regulator placed 103 of the 3\,249
broker-dealers it oversees in its proprietary-trading and market-making business segments; 93 of them had 150
registered representatives or fewer, and none had 500. Across all broker-dealers, almost half had ten registered people
or fewer, while the 149 largest employed 82\% of all registered representatives. The thousand-person firms of
chapters~2 to~4 are the tail of a long distribution whose body is firms of a few dozen people, many of them trading at
horizons of minutes to days. This chapter describes that body: what medium frequency means, how firm sizes are
distributed, how firms enter and leave, and the arrangements --- profit splits, arcades, funded-trader schemes --- that
a job seeker meets at the small end, some of them not jobs at all.

\section{Medium frequency: horizons, capacity and infrastructure}

\begin{definition}[Medium-frequency trading]\label{def:in:medium-frequency-proprietary-shops:mft}
\emph{Medium-frequency trading}\index{medium-frequency trading} is systematic trading whose positions last from minutes to
a few days: slower than high-frequency trading, whose edge is measured in microseconds and holding periods in seconds
(Book~10, chapter~24), and faster than the daily-to-monthly signals of most systematic funds.
\end{definition}

The horizon decides the firm. A medium-frequency strategy earns more per trade than a market maker and trades less
often, so its costs are market impact and fees rather than the price of being first; its infrastructure needs good data
and execution (Book~10) but not the tick-to-trade engineering of Book~13; and its capacity, the capital it can run before
its own trading erodes its edge (strategy capacity, Book~7, chapter~28), is larger than a market maker's and smaller than
a trend follower's. Its research is Book~7's and its strategies Book~8's: intraday and multi-day predictors on shares,
futures and currencies. A small team with good data and a disciplined research process can run such a book, which is
why the category contains so many small firms.

\begin{remark}[Proprietary, not necessarily small]\label{rem:in:medium-frequency-proprietary-shops:large}
Large firms trade at medium frequency too, often inside multi-strategy funds and platforms (chapters~4 and~5). The
chapter is about the standalone proprietary firms for which it is the whole business: firms that trade only their own
and their partners' capital.
\end{remark}

\section{The size distribution of proprietary firms}

The only public census of trading-firm sizes is the broker-dealer regulator's, and its unit is the registered
representative: a person registered to deal with the public, not an employee. A proprietary firm with no customers may
register few of its staff. The census still says how the industry's firms are shaped.

\begin{dated}{2024-12}{Broker-dealers by size}\label{dat:in:medium-frequency-proprietary-shops:sizes}
FINRA, end of 2024: 3\,249 registered broker-dealers. By registered representatives: 1\,522 firms with ten or fewer,
328 with 11--15, 205 with 16--20, 134 with 21--25, 108 with 26--30, 159 with 31--40, 102 with 41--50, 143 with 51--75, 80 with
76--100, 110 with 101--150, 137 with 151--300, 72 with 301--499, 65 with 500--1\,000 and 84 with more than 1\,000. Of the
649\,241 registered representatives, 81.7\% work at the 149 large firms (500 or more), 8.9\% at mid-size firms and
9.4\% at small ones. In its business segments, FINRA counts 93 small firms and 10 mid-size firms in proprietary trading
and market making, and none among the large.
\end{dated}

Most firms are small and most people work at large ones: 46.8\% of broker-dealers have ten registered representatives
or fewer, and 78.7\% have fifty or fewer, but those small firms employ 9.4\% of the representatives. A distribution
with that shape is usually described by its upper tail.

\begin{method}[Fitting a tail to binned counts]\label{met:in:medium-frequency-proprietary-shops:fit}
\begin{enumerate}
\item Choose a lower bound $x_{\min}$ and keep the bins at or above it.
\item Write the probability of each bin under a Pareto tail $P(X>x)=(x/x_{\min})^{-\alpha}$, treating integer bounds as
continuous at half-units; the log-likelihood is $\sum_i n_i\ln p_i(\alpha)$.
\item Maximise over $\alpha$; the standard error is the inverse square root of minus the second derivative of the
log-likelihood at the maximum.
\item Compare with another shape (a lognormal) by the chi-square distance between observed and expected counts, and
move $x_{\min}$ to see whether $\alpha$ is stable.
\end{enumerate}
\end{method}

From eleven registered representatives up (1\,727 firms, thirteen bins), the fitted tail index is $\alpha=0.60$ with a
standard error of 0.015 (\cref{fig:in:medium-frequency-proprietary-shops:ccdf}). An index below one is a heavy tail: in
such a distribution the largest firms hold a share of the total that does not shrink as firms are added, which is what
the representatives' split says directly. The fit is imperfect. The Pareto line overstates the share of the very
largest firms (6.5\% of firms above 1\,000 representatives against 4.9\% observed), a lognormal fits the bins better (a
chi-square distance of 5.1 against 21.2), and the index rises from 0.60 to 0.65 as the lower bound moves from 11 to 51.
The honest reading is a range, not a number: a tail index between about 0.6 and 0.65, a heavy tail by any of them.

\begin{omfigure}
\begin{tikzpicture}
\begin{axis}[width=10.5cm,height=6cm,xmode=log,ymode=log,xmin=9,xmax=1500,ymin=0.03,ymax=1.2,log ticks with fixed point,
  xlabel={\small registered representatives, $x$ (log scale)},ylabel={\small share of firms with at least $x$ (log scale)},
  tick label style={font=\footnotesize},grid=major,grid style={gray!20},table/col sep=comma,
  legend style={legend cell align=left,font=\scriptsize,at={(0.5,-0.3)},anchor=north,/tikz/every even column/.append style={column sep=8pt}},legend columns=2]
\addplot[only marks,mark=*,omDef] table[x=size,y=observed]{figdata/industry/06-medium-frequency-proprietary-shops/ccdf.csv};
\addplot[omThm,thick] table[x=size,y=pareto]{figdata/industry/06-medium-frequency-proprietary-shops/ccdf.csv};
\legend{observed (FINRA bins),{Pareto tail, $\alpha=0.60$}}
\end{axis}
\end{tikzpicture}

\omcaption{The share of the 1\,727 broker-dealers with eleven or more registered representatives that have at least $x$,
end of 2024, on log scales, with the fitted Pareto tail. Data: FINRA Industry Snapshot 2025, table 2.1.9, through
\texttt{in\_firmsize.fit}.}\label{fig:in:medium-frequency-proprietary-shops:ccdf}
\end{omfigure}

\section{Newer entrants and firms that grew out of crypto}

Firms enter and leave the registered population every year
(\cref{fig:in:medium-frequency-proprietary-shops:flows}). Between 2011 and 2024, on average 5.6\% of the previous year's
broker-dealers left FINRA membership and 3.2\% entered; in 2024, 185 left and 136 entered. The population
shrank by 29\% from 2010 to 2024 because exits outnumbered entries every year, not because entry stopped. Among
proprietary firms, a new entrant can be a handful of traders and researchers with their own capital or a backer's,
and an exit may be a firm closing, merging, or leaving the regulator's perimeter by trading only futures or only crypto
assets, which a broker-dealer registration does not cover.

The firms that grew out of crypto markets are that last case in reverse. The crypto trading firms of chapter~10 began
outside the securities regulators' registers, trading round the clock on venues with their own rules (Book~3, chapter~24),
and a crypto firm that extends to futures,
equities or options takes on the registrations that requires. For an employee
their appeal is the one of any young firm: a small team, wide responsibility, and pay that depends heavily on the
firm's own results.

\begin{omfigure}
\begin{tikzpicture}
\begin{axis}[width=11cm,height=5.2cm,xmin=2010.5,xmax=2024.5,ymin=0,ymax=8,ylabel={\small \% of previous year's firms},
  xtick={2011,2014,2017,2020,2023},xticklabel style={/pgf/number format/1000 sep={}},tick label style={font=\footnotesize},
  grid=major,grid style={gray!20},table/col sep=comma,
  legend style={legend cell align=left,font=\scriptsize,at={(0.5,-0.22)},anchor=north,/tikz/every even column/.append style={column sep=8pt}},legend columns=2]
\addplot[omThm,thick,mark=*] table[x=year,y=exit_pct]{figdata/industry/06-medium-frequency-proprietary-shops/flows.csv};
\addplot[omDef,thick,mark=square*] table[x=year,y=entry_pct]{figdata/industry/06-medium-frequency-proprietary-shops/flows.csv};
\legend{firms leaving,firms entering}
\end{axis}
\end{tikzpicture}

\omcaption{Broker-dealers leaving and entering FINRA membership each year, 2011--2024, in per cent of the previous
year-end population. Data: FINRA Industry Snapshot 2025, table 2.2.3, through
\texttt{in\_firmsize.entry\_exit}.}\label{fig:in:medium-frequency-proprietary-shops:flows}
\end{omfigure}

\section{Profit splits, arcades and funded-trader programmes}

At the small end the line between a job and a business blurs. Three arrangements recur.

\textbf{The profit split.} A trader at a small proprietary firm is often paid a share of the P\&L they generate, net of
their costs, with little or no salary: a formulaic payout in Book~16's terms (chapter~10), with the firm's capital,
technology and risk management as the firm's side of the bargain. The share, the costs charged and whether losses are
carried forward decide what the trader earns (chapter~13).

\begin{definition}[Trading arcade]\label{def:in:medium-frequency-proprietary-shops:arcade}
A \emph{trading arcade}\index{trading arcade} is a structure in which several legal entities or natural persons trade
under one firm's licence for proprietary trading, usually each with its own capital allocation, sharing the licence
holder's infrastructure, risk management and compliance.
\end{definition}

The Dutch regulator describes arcades in these words and states that they are possible in principle provided that all
the requirements on investment firms are met, with every participant's rights, powers and obligations clearly
allocated and documented. The capital rules that bind the licence holder (own funds and the K-factor requirement, Book~16,
chapter~2) bind the whole arcade: one participant's losses are the licence holder's problem.

\begin{definition}[Funded-trader programme]\label{def:in:medium-frequency-proprietary-shops:funded}
A \emph{funded-trader programme}\index{funded-trader programme} is a business that sells retail customers the chance to
trade the operator's capital: customers pay a fee to take an evaluation, usually on a simulated account, and those who
meet its targets are promised a share of the profits of a ``funded'' account.
\end{definition}

A \omterm{def:in:medium-frequency-proprietary-shops:funded}{funded-trader programme} is not an employer, and its customers are not traders of a proprietary firm: they pay the
operator, which earns its revenue from their fees. The public record shows how such a business can end up before a
regulator, and how such a case can end.

\begin{dated}{2025-05}{One funded-trader case in the court record}\label{dat:in:medium-frequency-proprietary-shops:mff}
In August 2023 the CFTC sued the operator of a \omterm{def:in:medium-frequency-proprietary-shops:funded}{funded-trader programme} in the US District Court for the District of New
Jersey, alleging that more than 135\,000 customers had signed up, paying at least \$310 million in fees; that some
account types required customers to generate simulated profits in a demonstration account first; and that funded accounts
were promised ``up to 85\% of profits''. On 13 May 2025 a special master appointed by the court, ruling on the
defendants' motion for sanctions against the CFTC, recommended that the motion be granted and the complaint be dismissed
with prejudice. The court dismissed the case with prejudice and awarded the defendants fees and costs of the sanctions
motion. The allegations were never proved.
\end{dated}

\section{What a small firm means for the employee}

A job at a small proprietary firm differs from one at a large firm in four ways that follow from the numbers above.
\begin{enumerate}
\item \textbf{The capital is the partners'.} Losses reduce what the owners have, and the firm's capital requirement
(Book~16, chapter~2) caps how much risk the whole firm can take; a bad month can shrink everyone's allocation.
\item \textbf{Pay follows P\&L closely.} Formulaic splits, few salaried roles, and deferred pay that is often absent.
\item \textbf{Controls are lighter.} Risk, compliance and operations may be one person each; the employee does more of
everything, including checking their own work.
\item \textbf{The firm may not last.} About one in eighteen broker-dealers leaves the register every year. The job's option value is what it teaches and who it introduces.
\end{enumerate}

\begin{example}[One trader's split]\label{ex:in:medium-frequency-proprietary-shops:split}
A trader at a small firm keeps 40\% of their book's net P\&L after costs of \$150\,000 a year, with no salary. A year of
\$800\,000 gross P\&L pays $0.4\times(800\,000-150\,000)=\$260\,000$; a year of \$200\,000 pays \$20\,000; a year of
\$100\,000 pays nothing, and if losses are carried forward the next year's first \$50\,000 of net P\&L pays nothing either.
The same trader at a salaried firm would earn a base salary and a discretionary bonus with a lower variance and, usually,
a lower expected total (chapter~13 values such offers).
\end{example}

\section{Tutorial: the long tail}\label{tut:in:medium-frequency-proprietary-shops:tut}

\begin{tutorial}
\textbf{Goal.} Fit the size distribution of broker-dealers from FINRA's published bins and read entry and exit.
\textbf{End state:} \cref{fig:in:medium-frequency-proprietary-shops:ccdf} and the tail index with its standard error.
\begin{tutsteps}
\item \textbf{The bins.} \texttt{data/industry/finra\_size\_bins.csv} holds FINRA's fourteen size bins for 2020--2024;
\texttt{in\_firmsize.bins(2024)} reads them as \texttt{firm.firmsize.Bin} records.
\item \textbf{Fit.} \texttt{firm.firmsize.pareto\_fit(bins, 11)} maximises the binned likelihood
(\cref{lst:in:medium-frequency-proprietary-shops:fit}); \texttt{lognormal\_fit} fits the comparison.
\omcode{code/firm/firmsize/firm_firmsize.py}{51}{58}{The binned Pareto fit and its standard error from the observed
information.}\label{lst:in:medium-frequency-proprietary-shops:fit}
\item \textbf{Check.} \texttt{expected\_counts} and \texttt{chi2} compare the two fits; refit with $x_{\min}$ of 16, 21,
31 and 51 and watch $\alpha$ move from 0.60 to 0.65.
\item \textbf{Flows.} \texttt{in\_firmsize.entry\_exit()} computes the exit and entry rates behind
\cref{fig:in:medium-frequency-proprietary-shops:flows}.
\end{tutsteps}
\end{tutorial}

\textbf{What to change next.} Fit the 2020 bins and compare the index with 2024's (exercise~7); simulate firm sizes from
the fitted tail with \texttt{sample\_pareto} and check that the fit recovers the index you drew from.

\section{Build: size distributions from binned counts}\label{bld:in:medium-frequency-proprietary-shops:firmsize}

\begin{build}
\bfield{Purpose} Fit size distributions to the binned counts regulators publish, for firm sizes here and for pay bands
and headcounts later (chapters~12 and~14).
\bfield{Interface} \texttt{firm.firmsize}: \texttt{Bin(lo, hi, n)}; \texttt{pareto\_fit(bins, x\_min)};
\texttt{lognormal\_fit(bins, x\_min)}; \texttt{expected\_counts}; \texttt{chi2}; \texttt{ccdf}; \texttt{sample\_pareto}.
\bfield{Rules} Integer bounds are continuous at half-units; an open top bin has probability equal to the tail beyond its
lower bound; a fit states its lower bound and its standard error.
\bfield{Acceptance tests} \texttt{code/firm/firmsize/tests/}: the fit recovers a known index from 20\,000 simulated
sizes within three standard errors; expected counts sum to the observed total; the tail share falls with size; the
lognormal fit runs.
\bfield{Stretch} Choose $x_{\min}$ by minimising the distance between the fitted and observed tail (the Clauset--Shalizi--Newman
method); a bootstrap interval for $\alpha$ over bins.
\end{build}

\begin{omsources}
\item FINRA, Industry Snapshot 2025: tables 1.1.3, 2.1.9, 2.2.3 and 2.6.1--2.6.3.
\item AFM (Netherlands Authority for the Financial Markets), Market participants subject to a licence obligation, on
trading arcades.
\item CFTC v.\ Traders Global Group Inc.\ et al., D.N.J.\ No.~23-cv-11808: complaint (29 August 2023); report and
recommendation of the special master (13 May 2025); Willkie Farr \& Gallagher's summary of the court's order (26 May 2025).
\item A.~Clauset, C.~R.~Shalizi and M.~E.~J.~Newman, ``Power-law distributions in empirical data'', \emph{SIAM Review}
51(4), 2009.
\end{omsources}

\section{Exercises}

\begin{exercise}[$\star$]\label{exo:in:medium-frequency-proprietary-shops:1}
From the dated box, what share of broker-dealers had more than 500 registered representatives at the end of 2024, and
what share of all representatives did they employ?
\end{exercise}

\begin{exercise}[$\star$]\label{exo:in:medium-frequency-proprietary-shops:2}
In the example, what does the trader earn on gross P\&L of \$500\,000, and on \$150\,000?
\end{exercise}

\begin{exercise}[$\star$]\label{exo:in:medium-frequency-proprietary-shops:3}
Under a Pareto tail with $\alpha=0.6$ from $x_{\min}=10.5$, what is the probability that a firm with at least eleven
representatives has more than 1\,000?
\end{exercise}

\begin{exercise}[$\star\star$]\label{exo:in:medium-frequency-proprietary-shops:4}
Compute the average exit and entry rates of broker-dealers over 2011--2024 and the rate of net shrinkage they imply.
\end{exercise}

\begin{exercise}[$\star\star$]\label{exo:in:medium-frequency-proprietary-shops:5}
Why is a tail index below one described as heavy? What does it imply for the share of the largest firms as the
population grows?
\end{exercise}

\begin{exercise}[$\star\star$]\label{exo:in:medium-frequency-proprietary-shops:6}
From \cref{fig:in:medium-frequency-proprietary-shops:ccdf}, where does the Pareto line lie above the observed shares,
and what does that say about the largest firms?
\end{exercise}

\begin{exercise}[$\star\star\star$]\label{exo:in:medium-frequency-proprietary-shops:7}
\emph{Coding.} Fit the 2020 bins from $x_{\min}=11$. Is the index different from 2024's by more than two standard errors?
\end{exercise}

\begin{exercise}[$\star\star\star$]\label{exo:in:medium-frequency-proprietary-shops:8}
\emph{Find the flaw.} ``Half of all broker-dealers have ten staff or fewer, so a typical trading job is at a firm of
ten.''
\end{exercise}

\section{Problem: The Long Tail}

\begin{problem}[{Weekend problem --- the long tail}]\label{pb:in:medium-frequency-proprietary-shops:1}
A researcher with an offer from a twelve-person medium-frequency firm wants to understand where such firms sit in the
industry and what the arrangement means.

\medskip\noindent\textbf{Part I --- The category.}
\begin{enumerate}
\item Define \omterm{def:in:medium-frequency-proprietary-shops:mft}{medium-frequency trading} and place it between high-frequency trading and a systematic fund.
\item Why do its costs and infrastructure differ from a market maker's?
\item What unit does FINRA use for firm size, and why does it understate a proprietary firm?
\item How many broker-dealers did FINRA place in its proprietary-trading and market-making segments, and how large were they?
\item What share of broker-dealers had ten registered representatives or fewer, and fifty or fewer?
\end{enumerate}

\medskip\noindent\textbf{Part II --- The tail.}
\begin{enumerate}[resume]
\item Write the probability of one bin under a Pareto tail with half-unit bounds.
\item Give the fitted index from eleven representatives and its standard error.
\item Compare the Pareto and lognormal fits by their chi-square distances.
\item How does the index move with the lower bound?
\item What do the representatives' shares by firm size say that the index also says?
\end{enumerate}

\medskip\noindent\textbf{Part III --- Entry, exit and arrangements.}
\begin{enumerate}[resume]
\item Give the average exit and entry rates over 2011--2024, and the 2024 counts.
\item Define a \omterm{def:in:medium-frequency-proprietary-shops:arcade}{trading arcade}, and say what the Dutch regulator requires of one.
\item Define a \omterm{def:in:medium-frequency-proprietary-shops:funded}{funded-trader programme} and say why it is not an employer.
\item State the funded-trader case's allegations, its outcome and its date, neutrally.
\item Compute the trader's pay in the example for gross P\&L of \$800\,000, \$200\,000 and \$100\,000.
\end{enumerate}

\medskip\noindent\textbf{Part IV --- The verdict.}
\begin{enumerate}[resume]
\item State the \emph{named result}: the tail index with its standard error and its range as the lower bound moves, and
the share of representatives at small firms.
\item Name the four ways a job at a small firm differs from one at a large firm.
\item What should the researcher ask about capital, pay and controls?
\item Roughly what is the chance that a broker-dealer leaves the register in a given year?
\item In two sentences, how should the researcher weigh the offer?
\end{enumerate}
\end{problem}

\section{Interview questions}

\begin{interviewq}[$\star$ \iqroles{researcher, trader}]\label{iq:in:medium-frequency-proprietary-shops:1}
What changes in research and infrastructure when a strategy's holding period goes from seconds to a day?
\end{interviewq}

\begin{interviewq}[$\star$ \iqroles{trader}]\label{iq:in:medium-frequency-proprietary-shops:2}
You are offered 50\% of your net P\&L with no salary, or a salary of \$150\,000 and a 10\% bonus of net P\&L. At what
net P\&L are they equal?
\end{interviewq}

\begin{interviewq}[$\star\star$ \iqroles{researcher}]\label{iq:in:medium-frequency-proprietary-shops:3}
You have only binned counts of firm sizes. How do you estimate a tail index, and how do you know whether a power law is
the right shape?
\end{interviewq}

\begin{interviewq}[$\star\star$ \iqroles{risk}]\label{iq:in:medium-frequency-proprietary-shops:4}
Several traders trade under one firm's licence, each with their own capital. What risks does the licence holder carry?
\end{interviewq}

\begin{interviewq}[$\star\star$ \iqroles{researcher}]\label{iq:in:medium-frequency-proprietary-shops:5}
A medium-frequency strategy makes 5 basis points a trade before costs, trades each position once a day on
\$20 million of capital with a turnover of 100\% a day, and pays 2 basis points in costs. What is its annual net return
on capital, and what limits scaling it?
\end{interviewq}

\begin{interviewq}[$\star\star\star$ \iqroles{researcher, risk}]\label{iq:in:medium-frequency-proprietary-shops:6}
A business sells evaluations to would-be traders and promises profit splits to those who pass. How would you tell from
its numbers whether its revenue comes from trading or from fees?
\end{interviewq}
